% Flesch Reading Ease on NeurIPS abstracts (body figure, NeurIPS only).
% Cross-venue and all-15-metric versions are in the appendix.
% Data: ../data/pgfplots/readability_flesch_ease.csv
% Higher values = easier reading.

\begin{figure}[t]
\centering
\begin{tikzpicture}
\begin{axis}[
    width=0.92\linewidth, height=5cm,
    xlabel={Year},
    ylabel={Flesch Reading Ease (\textuparrow\ easier)},
    xtick distance=4,
    x tick label style={rotate=45, anchor=east, font=\tiny,
                        /pgf/number format/1000 sep={}},
    ytick align=inside,
    tick label style={font=\tiny},
    label style={font=\small},
    enlarge x limits=0.02,
  ]
  \addplot[color={rgb,255:red,44;green,160;blue,44}, solid, very thick]
    table[x=year, y=neurips, col sep=comma]
    {../data/pgfplots/readability_flesch_ease.csv};
  % ChatGPT launch marker (late 2022)
  \draw[black, dashed, semithick] (axis cs:2022,-5) -- (axis cs:2022,50);
  \node[font=\tiny, rotate=90, anchor=south west]
    at (axis cs:2022.2,18) {ChatGPT (2022)};
\end{axis}
\end{tikzpicture}
\caption{\textbf{Flesch Reading Ease on NeurIPS abstracts,
  1987--2025.}
  Higher values indicate easier reading (\textuparrow). NeurIPS
  falls from approximately 24 in 1987 to approximately 8 in 2025,
  a sustained decline. All 15 classical metrics move in the same
  direction on NeurIPS (Appendix~\ref{app:readability}). Vertical dashed line marks the
  late-2022 availability of widely-used instruction-tuned writing
  assistants.}
\label{fig:readability-flesch}
\end{figure}
